% =============================================================================
% docs/theory.tex  ---  T18
% Edited under T25: attributions that could not be opened removed, and the
% elementary bounds they supported are now proved in the text.
% Fragment for the manuscript. No preamble; include inside the body with \input.
% Requires: amsmath, amsthm.
%
% The five \newtheorem lines below are part of the fragment so that it compiles
% on its own under a bare \documentclass{article} wrapper.  If the manuscript
% preamble already declares these environments, delete this block instead.
% =============================================================================
% Theorem environments are declared once, in the manuscript preamble.
% For a standalone compile, declare them in the wrapper.

% =============================================================================
\subsection{The count functional: what the marginals determine}
% =============================================================================

Let the country be divided into $n$ areas, $s=1,\dots,n$.  Area $s$ has latent
prevalence $\theta_s$, and for a crisis line $c$ define the indicator
$I_s = \mathbf{1}\{\theta_s > c\}$ and the marginal exceedance probability
$p_s = \operatorname{E}[I_s] = \operatorname{P}(\theta_s > c)$.  Write
$P_{st} = \operatorname{P}(I_s = 1,\ I_t = 1)$ for $s \neq t$,
$S = \sum_{s=1}^{n} p_s$, and
\begin{equation}
  C \;=\; \sum_{s=1}^{n} I_s \;=\; \#\{\, s : \theta_s > c \,\},
  \label{eq:count}
\end{equation}
the number of areas above the line.  $\operatorname{Var}(C)$ refers to the
variance under a \emph{joint} distribution for $(\theta_1,\dots,\theta_n)$;
the marginals $(p_s)$ are held fixed throughout.  A single number
$\operatorname{E}[C]$ and a single interval are reported in the application,
but the propositions below say exactly which of the reported quantities the
marginals can carry.

% -----------------------------------------------------------------------------
\begin{proposition}[mean and variance of the count]
\label{prop:meanvar}
For any joint distribution of $(\theta_1,\dots,\theta_n)$ with
$\operatorname{P}(\theta_s > c) = p_s$,
\begin{align}
  \operatorname{E}[C] &= \sum_{s=1}^{n} p_s, \label{eq:mean}\\
  \operatorname{Var}(C) &= \sum_{s=1}^{n} p_s(1-p_s)
     \;+\; 2\sum_{1 \le s < t \le n}\bigl(P_{st} - p_s p_t\bigr). \label{eq:var}
\end{align}
The right-hand side of \eqref{eq:mean} involves the marginals only.  The second
sum in \eqref{eq:var} vanishes for \emph{every} joint distribution with these
marginals if and only if $P_{st} = p_s p_t$ for all $s<t$.
\end{proposition}

\begin{proof}
$C$ is a finite sum and expectation is linear, so
$\operatorname{E}[C] = \sum_s \operatorname{E}[I_s] = \sum_s p_s$, which is
\eqref{eq:mean}; no property of the joint distribution is used.  For
\eqref{eq:var}, $\operatorname{Var}(C) = \sum_s \operatorname{Var}(I_s)
+ 2\sum_{s<t}\operatorname{Cov}(I_s,I_t)$.  Since $I_s$ is Bernoulli with mean
$p_s$, $\operatorname{Var}(I_s) = p_s(1-p_s)$, and
$\operatorname{Cov}(I_s,I_t) = \operatorname{E}[I_s I_t] - p_s p_t = P_{st} - p_s p_t$.

For the last sentence: if $P_{st} = p_s p_t$ for all $s<t$ then every summand
is zero.  Conversely, if some $P_{s_0 t_0} \ne p_{s_0} p_{t_0}$, choose the joint
distribution that differs from the product law only in the
$(s_0,t_0)$ cell; then the second sum is $2(P_{s_0t_0} - p_{s_0}p_{t_0}) \ne 0$.
Hence the sum vanishes identically on the set of joints with these marginals
only under pairwise independence.
\end{proof}

\begin{remark}[the converse fails, and this matters for exposition]
\label{rem:converse}
For a \emph{particular} joint distribution the second sum can vanish while the
indicators remain dependent, because covariances of opposite sign cancel.  With
$n=3$, $p_1=p_2=p_3=\tfrac12$ and the law
\[
  \begin{array}{c|cccccccc}
    (I_1,I_2,I_3) & (1,1,1) & (1,1,0) & (1,0,1) & (0,1,1) & (1,0,0) & (0,1,0) & (0,0,1) & (0,0,0)\\[2pt]
    \hline
    \text{probability} & 0.10 & 0.25 & 0.15 & 0.05 & 0.00 & 0.10 & 0.20 & 0.15
  \end{array}
\]
the marginals are all $\tfrac12$, the covariances are $+0.1$, $0$ and $-0.1$,
the variance is $\operatorname{Var}(C) = 0.75$, exactly the value
$\sum_s p_s(1-p_s)$ that independence would give, and yet
$\operatorname{P}(I_1=1,I_3=1)=0.25 \neq p_1p_3$.  The correct statement is
one-directional: pairwise uncorrelatedness implies the second sum vanishes;
a vanishing second sum does not imply pairwise uncorrelatedness for $n \ge 3$.
(The equivalence does hold when $n=2$.)
\end{remark}

\begin{remark}[what a marginal-only report can and cannot claim]
\label{rem:consequence}
By Proposition~\ref{prop:meanvar}, $\operatorname{E}[C]$ is a function of the
marginals alone: any method that publishes per-area probabilities $p_s$
reproduces $\operatorname{E}[C]$ exactly, however it obtained them.  Every
other functional of $C$ --- its variance, its quantiles, the probability that
the count exceeds a policy-relevant number --- depends on the joint
distribution, and Proposition~\ref{prop:sharp} shows that the marginals leave
that joint free over a range wide enough to reverse qualitative conclusions.
A marginal-only method is therefore exact about the mean and unevidenced about
everything else.
\end{remark}

% -----------------------------------------------------------------------------
\begin{proposition}[sharp bounds on the variance given the marginals]
\label{prop:sharp}
Fix $p_1,\dots,p_n \in [0,1]$, let $S=\sum_s p_s$, $k=\lfloor S \rfloor$ and
$f = S-k \in [0,1)$.  Then
\begin{equation}
  f(1-f) \;\le\; \operatorname{Var}(C) \;\le\;
  \sum_{s=1}^{n}\sum_{t=1}^{n} \min(p_s,p_t) \;-\; S^2 ,
  \label{eq:sharp}
\end{equation}
both endpoints are attained by explicit joint distributions with these
marginals, and every value in the interval is attained.
\end{proposition}

\begin{proof}
\emph{Upper end.}  For each pair, $\{I_s=1,\,I_t=1\}$ is a subset of
$\{I_s=1\}$ and of $\{I_t=1\}$, so $P_{st} \le \min(p_s,p_t)$; that pairwise
inequality is all the upper end uses.  The right-hand side of
\eqref{eq:var} is increasing in each $P_{st}$ with coefficient $2>0$, so
\[
  \operatorname{Var}(C) \le \sum_s p_s(1-p_s)
   + 2\sum_{s<t}\bigl(\min(p_s,p_t) - p_s p_t\bigr)
   = \sum_{s,t}\min(p_s,p_t) - S^2 ,
\]
using $\sum_s p_s(1-p_s) = S - \sum_s p_s^2$ and
$S^2 = \sum_s p_s^2 + 2\sum_{s<t}p_sp_t$.  Attainment is Lemma~\ref{lem:comon}.

\emph{Lower end.}  The bound uses only that $C$ is integer-valued with mean
$S$; it is Lemma~\ref{lem:floor}.  Attainment is Lemma~\ref{lem:cyclic}.

\emph{Every intermediate value.}  The set of joint distributions with fixed
marginals is convex and $\operatorname{Var}(C) = \operatorname{E}[C^2]-S^2$ is
linear in the joint law, so its image is an interval.
\end{proof}

\begin{lemma}[integer floor]
\label{lem:floor}
If $X$ is integer-valued with $\operatorname{E}[X]=S$, then
$\operatorname{Var}(X) \ge f(1-f)$ with $f = S-\lfloor S\rfloor$.
\end{lemma}

\begin{proof}
For every integer $x$ the parabola $(x-k)(x-k-1) \ge 0$ gives
$x^2 \ge (2k+1)x - k(k+1)$, with equality at $x=k$ and $x=k+1$.  Taking
expectations and using $\operatorname{E}[X]=k+f$,
\[
  \operatorname{E}[X^2] \;\ge\; (2k+1)(k+f) - k(k+1)
  \;=\; k^2 + 2kf + f .
\]
Since $k^2(1-f)+(k+1)^2 f = k^2+2kf+f$, this is
$\operatorname{E}[X^2] \ge k^2(1-f)+(k+1)^2f$, and subtracting
$S^2=(k+f)^2$ gives $\operatorname{Var}(X) \ge f(1-f)$.
\end{proof}

\begin{lemma}[comonotone coupling attains the upper end]
\label{lem:comon}
Let $U$ be uniform on $[0,1]$ and set $I_s = \mathbf{1}\{U \le p_s\}$.  Then
$\operatorname{P}(I_s=1)=p_s$ and
$\operatorname{Var}(C) = \sum_{s,t}\min(p_s,p_t) - S^2$.
\end{lemma}

\begin{proof}
$\operatorname{P}(U \le p_s) = p_s$, and
$\operatorname{P}(I_s=I_t=1) = \operatorname{P}(U \le \min(p_s,p_t)) =
\min(p_s,p_t)$, so every pairwise upper bound in the proof of
Proposition~\ref{prop:sharp} holds with equality simultaneously.  Substituting
into \eqref{eq:var} gives the stated value.
\end{proof}

\begin{lemma}[cyclic coupling attains the lower end]
\label{lem:cyclic}
Let $U$ be uniform on $[0,1]$, let $c_0=0$ and $c_s = p_1+\dots+p_s$, and set
\begin{equation}
  I_s \;=\; \mathbf{1}\bigl\{\, (U - c_{s-1}) \bmod 1 \;<\; p_s \,\bigr\},
  \qquad s = 1,\dots,n .
  \label{eq:cyclic}
\end{equation}
Then $\operatorname{P}(I_s=1)=p_s$ for every $s$, and $C$ takes only the two
values $k$ and $k+1$, with $\operatorname{P}(C=k+1)=f$.  Hence
$\operatorname{Var}(C)=f(1-f)$.
\end{lemma}

\begin{proof}
Each $I_s$ is an arc of the circle of circumference $1$ of length $p_s$, so
$\operatorname{P}(I_s=1)=p_s$.  For $u \in [0,1)$, $I_s(u)=1$ if and only if
$u+j \in [c_{s-1},c_s)$ for some integer $j$.  The intervals
$[c_{s-1},c_s)$, $s=1,\dots,n$, partition $[0,S)$; a point $u+j$ therefore lies
in exactly one of them when $u+j \in [0,S)$ and in none otherwise.  Hence
\[
  C(u) \;=\; \#\{\, j \text{ integer} : 0 \le u+j < S \,\}
        \;=\; \#\{\, j \text{ integer} : -u \le j < S-u \,\}.
\]
For $u \in (0,1)$ the bounds $-\lceil u \rceil = 0$ and $\lfloor S-u \rfloor$
give $C(u) = \lfloor k+f-u \rfloor + 1$, which equals $k+1$ when $u<f$ and $k$
when $u>f$.  Thus $\operatorname{P}(C=k+1)=f$ and the variance is
$f(1-f)$.
\end{proof}

\begin{corollary}[equal marginals, and the degenerate case]
\label{cor:equal}
If $p_s \equiv p$ then the upper end of \eqref{eq:sharp} is $n^2p(1-p)$, so
\[
  \frac{\operatorname{sd}(C)}{\sqrt{\sum_s p_s(1-p_s)}} \;\le\; \sqrt{n}
\]
with equality at the comonotone coupling; the ratio is $\sqrt n$, not $n$.  If
$np$ is an integer then $f=0$, the lower end is $0$, and
$C \equiv np$ under the coupling \eqref{eq:cyclic}: the count is then
determined by the marginals, and the ``lower bound'' is degenerate rather than
positive.  For the unequal marginals of the application the extremal ratio is
smaller than $\sqrt n$, because heterogeneity in $p_s$ shrinks
$\sum_{s,t}\min(p_s,p_t)$.
\end{corollary}

\begin{remark}[the pairwise route to the lower end is invalid]
\label{rem:pairwise}
It is tempting to insert the \emph{lower} bounds
$P_{st} \ge \max(0,p_s+p_t-1)$ into \eqref{eq:var} and call the result the
minimum.  Each of those bounds is one line: $P_{st} \ge 0$ is clear, and
inclusion--exclusion gives
$P_{st} = p_s+p_t-\operatorname{P}(I_s=1 \text{ or } I_t=1) \ge p_s+p_t-1$,
because a probability is at most $1$.  (The upper and lower pairwise bounds
together are commonly named the Fr\'echet--Hoeffding bounds.)  Inserting them
is nevertheless wrong for $n \ge 3$: the pairwise lower bounds need not be
attained by any single joint distribution.  For $p_s \equiv \tfrac12$ and $n=3$ all three
bounds are $0$, which would require $C \in \{0,1\}$ and hence $S \le 1$,
contradicting $S=\tfrac32$.  On the $774$ marginals of the application the
formula returns a \emph{negative} variance, and is therefore visibly
impossible.  The bound in \eqref{eq:sharp} must be used instead; it is
attained, and it is attained by \eqref{eq:cyclic}, not by a coupling that
satisfies every pairwise bound.
\end{remark}

\begin{remark}[what the marginals do not determine: the upper tail]
\label{rem:tail}
Two joint distributions with identical marginals give
$\operatorname{P}(C \ge 101) = 0.476$ (comonotone) and
$\operatorname{P}(C \ge 101) = 1.000$ (cyclic) on the application's own
$774$ marginals, against $0.984$ under the fitted model.  The extremal-variance
couplings need not extremise the tail probability, so these two numbers are a
demonstration that the tail is not a marginal quantity, not the endpoints of
its feasible range.
\end{remark}

% -----------------------------------------------------------------------------
\begin{proposition}[coverage of the identified set when the bounds are estimated]
\label{prop:coverage}
Let $T$ be a random variable, let $L \le U$ be functions of $T$, and let
$q_\beta(X) = \inf\{x : \operatorname{P}(X \le x) \ge \beta\}$.  Then for
$0<\alpha<1$,
\begin{equation}
  \operatorname{P}\Bigl(\,[L(T),U(T)] \subseteq
    \bigl[\,q_{\alpha/2}(L(T)),\ q_{1-\alpha/2}(U(T))\,\bigr]\,\Bigr)
  \;\ge\; 1-\alpha ,
  \label{eq:coverage}
\end{equation}
with no monotonicity assumption on $L$ or $U$.  If in addition $L$ and $U$ are
non-decreasing in $T$, then the event in \eqref{eq:coverage} is a single
interval of values of $T$; if they are continuous and strictly increasing the
event is exactly the central $(1-\alpha)$ interval of $T$ and
\eqref{eq:coverage} holds with equality.
\end{proposition}

\begin{proof}
For any real random variable $X$ and any $\beta \in (0,1)$, the definition of
$q_\beta$ as an infimum gives $\operatorname{P}(X < q_\beta(X)) \le \beta$, and
right-continuity of $x \mapsto \operatorname{P}(X \le x)$ gives
$\operatorname{P}(X \le q_\beta(X)) \ge \beta$.  Apply this to $L(T)$ with
$\beta=\alpha/2$ and to $U(T)$ with $\beta=1-\alpha/2$: the events
$A=\{L(T) \ge q_{\alpha/2}(L(T))\}$ and
$B=\{U(T) \le q_{1-\alpha/2}(U(T))\}$ satisfy
$\operatorname{P}(A) \ge 1-\alpha/2$ and
$\operatorname{P}(B) \ge 1-\alpha/2$.  Since $L \le U$, containment of the
identified set is exactly $A \cap B$, so
\[
  \operatorname{P}(A \cap B) \;\ge\; \operatorname{P}(A) + \operatorname{P}(B) - 1
  \;\ge\; 1-\alpha .
\]
For the refinement: if $L$ is non-decreasing then $A=\{T \ge a\}$ with
$a=\inf\{t : L(t) \ge q_{\alpha/2}(L(T))\}$, and if $U$ is non-decreasing then
$B=\{T \le b\}$ with $b=\sup\{t: U(t) \le q_{1-\alpha/2}(U(T))\}$, so
$A\cap B = \{a \le T \le b\}$ is an interval.  If both bounds are continuous
and strictly increasing, $a$ and $b$ are the $\alpha/2$ and $1-\alpha/2$
quantiles of $T$ itself, and the two tail probabilities add to exactly
$\alpha$.
\end{proof}

\begin{remark}[what Proposition~\ref{prop:coverage} does not give]
\label{rem:p3limits}
The probability in \eqref{eq:coverage} is taken over the distribution of $T$,
which in the application is the posterior distribution of the model.  It is a
posterior statement: it carries no frequentist coverage guarantee for the
identified set, and it conditions on the model for $\theta$ in the observed
areas being correct --- which is the object under dispute.  It treats the
observed/unobserved partition and the population share $p$ as known, so error
in assigning clusters to areas is outside it.  It is also a statement about one
estimand at a time: reporting a count and a population share together, each
with its own $95\%$ interval, gives no simultaneous guarantee for the pair.
Finally, the monotonicity hypothesis is what licenses the sentence ``the
identified set lies inside the reported interval exactly when $T$ lies in its
own central interval''.  That sentence is false without monotonicity even
though the inequality \eqref{eq:coverage} survives, so the hypothesis should be
stated and checked rather than assumed.
\end{remark}

% -----------------------------------------------------------------------------
% -----------------------------------------------------------------------------
\begin{lemma}[identified set when some areas are unobserved]
\label{lem:identified}
Partition the $n$ areas into an observed set $O$ and an unobserved set $U$,
with $w=|U|$, and suppose the observed statuses $I_s$, $s\in O$, are known.
Let $\omega_s$ be the population of area $s$, let $p$ be the population share
held by observed areas, let $C_O = \sum_{s\in O} I_s$, and let $\pi_O$ be the
population-weighted share above the line among observed areas:
\[
  p \;=\; \frac{\sum_{s\in O}\omega_s}{\sum_{s=1}^{n}\omega_s},
  \qquad
  \pi_O \;=\; \frac{\sum_{s\in O}\omega_s I_s}{\sum_{s\in O}\omega_s}.
\]
Then the identified set for the count is
\begin{equation}
  \bigl[\,C_O,\; C_O + w\,\bigr],
  \label{eq:identset-count}
\end{equation}
and the identified set for the population share above the line is
\begin{equation}
  \bigl[\,p\,\pi_O,\; p\,\pi_O + (1-p)\,\bigr].
  \label{eq:identset-share}
\end{equation}
Every value in each interval is attained by some configuration of the
unobserved areas.  In the application the observed statuses are read from a
posterior draw, so the two sets are conditional on that draw.
\end{lemma}

\begin{proof}
Write $C = C_O + C_U$ with $C_U = \sum_{s\in U} I_s$.  Each term of $C_U$ is
$0$ or $1$, so $0 \le C_U \le w$, and every integer in $[0,w]$ is attained by
placing exactly that many unobserved areas above $c$ and the rest below; no
assumption is imposed on $(\theta_s)_{s\in U}$ that would exclude such a
configuration.  For the share, let
$N = \sum_{s=1}^{n}\omega_s I_s / \sum_{s=1}^{n}\omega_s$ be the national
quantity and let $\pi_U = \sum_{s\in U}\omega_s I_s / \sum_{s\in U}\omega_s$
be the unobserved areas' share.  Then
\[
  N \;=\; p\,\pi_O + (1-p)\,\pi_U ,
\]
with $\pi_U \in [0,1]$.  The right-hand side is increasing in $\pi_U$, so its
image over $[0,1]$ is \eqref{eq:identset-share}, and $\pi_U=0$ and $\pi_U=1$
are both attained.
\end{proof}

% -----------------------------------------------------------------------------
\begin{proposition}[monotonicity in the crisis line]
\label{prop:monotone}
For $c < c'$, $C(c') \le C(c)$ for every realisation of
$(\theta_1,\dots,\theta_n)$.  Consequently $\operatorname{E}[C(c)]$ is
non-increasing in $c$, the identified sets are nested,
$[L(c'),U(c')] \subseteq [L(c),U(c)]$, and both endpoints of the reported
interval --- and every quantile of $C(c)$ --- are non-increasing in $c$.
\end{proposition}

\begin{proof}
If $\theta_s > c'$ then $\theta_s > c$, so
$\mathbf{1}\{\theta_s>c'\} \le \mathbf{1}\{\theta_s>c\}$ for every $s$ and every
outcome; summing over $s$ gives $C(c') \le C(c)$ pathwise.  Expectation and
quantiles are monotone with respect to pathwise domination, so
$\operatorname{E}[C(c')] \le \operatorname{E}[C(c)]$ and
$q_\beta(C(c')) \le q_\beta(C(c))$ for every $\beta$.  For nesting,
Lemma~\ref{lem:identified} gives the identified set $[C_O(c), C_O(c)+w]$
for the count and $[p\,\pi_O(c), p\,\pi_O(c)+(1-p)]$ for the population
share.  Neither width depends on $c$, and $C_O(c') \le C_O(c)$ and
$\pi_O(c') \le \pi_O(c)$ pathwise, so
$[C_O(c'), C_O(c')+w] \subseteq [C_O(c), C_O(c)+w]$, and likewise for the
share.
\end{proof}

\begin{remark}
\label{rem:sdnot}
Proposition~\ref{prop:monotone} does not extend to the standard deviation:
a stochastically non-increasing family can have non-monotone spread.  In the
application $\operatorname{sd}(C(c))$ does decrease over
$c \in \{0.4,0.5,0.6\}$, but that is an observation on one fit, not a
consequence of the proposition, and it should not be presented as one.
\end{remark}

% -----------------------------------------------------------------------------
\subsection*{Numerical values for this application}

Table~\ref{tab:countnumbers} reports each quantity in
Propositions~\ref{prop:meanvar}--\ref{prop:monotone} evaluated on the $774$
area-level exceedance probabilities of the fitted model.  Every value is
produced by \texttt{theory\_counts.R}, which also simulates the two
extremal couplings and reproduces the closed forms.

\begin{table}[htbp]
\centering
\begin{tabular}{lr}
\hline
quantity & value \\
\hline
areas $n$ & $774$ \\
$\sum_s p_s = \operatorname{E}[C]$ & $114.975$ \\
$\sum_s p_s(1-p_s)$ (independence) & $40.258$ \\
lower end $f(1-f)$ of \eqref{eq:sharp} & $0.0244$ \\
upper end $\sum_{s,t}\min(p_s,p_t)-S^2$ & $5492.63$ \\
fitted posterior variance & $54.509$ \\
$\operatorname{sd}(C)$ at lower end / fitted / upper end & $0.156$ / $7.383$ / $74.112$ \\
extremal ratio $R$ at lower end / fitted / upper end & $0.025$ / $1.164$ / $11.681$ \\
\hline
\end{tabular}
\caption{The count functional at $c=0.50$ on $n=774$ areas.  The mean is a
marginal quantity; the variance is not.  The fitted variance lies at $0.99\%$
of the way up its feasible range, and the marginals alone permit a standard
deviation forty-seven times smaller or ten times larger than the fitted one.}
\label{tab:countnumbers}
\end{table}

% -----------------------------------------------------------------------------
\subsection*{Attribution}

Nothing above is new method.  The ingredients are classical, and
Table~\ref{tab:attribution} records what each one is and where it is dealt
with; what is assembled here is the statement of what the marginals determine
for this count, and its consequence for a method that reports per-area
probabilities.

Two of those ingredients carry names that readers will recognise, and neither
is given a bibliography entry.  The pairwise bounds of
Proposition~\ref{prop:sharp} and Remark~\ref{rem:pairwise} are commonly called
the Fr\'echet--Hoeffding bounds, and the identified sets of
Lemma~\ref{lem:identified} are commonly called Manski bounds.  Both are
elementary, and both are proved in the text where they are used --- two lines
and one line respectively --- so neither rests on a citation.  The names are
kept for orientation only.

One result here is deliberately not claimed, and one citation is deliberately indirect. The interval
\eqref{eq:coverage} is \emph{not} a confidence set for a partially identified parameter in the sense of
Imbens and Manski (2004). Their object is built to cover the true value of the parameter with a fixed
probability; ours is built to contain the whole identified set, which is the more demanding requirement and
is what the union bound buys us. The two are routinely confused in applied work, which is why the distinction
is stated rather than left to the reader. Selection among interval estimators for partially identified
parameters is theirs, not ours, and we take no position on it.

\begin{table}[htbp]
\centering
\begin{tabular}{p{0.30\textwidth}p{0.28\textwidth}p{0.32\textwidth}}
\hline
result & classical source & status here \\
\hline
pairwise bounds $P_{st} \le \min(p_s,p_t)$ and $P_{st} \ge \max(0,p_s+p_t-1)$ & commonly named the Fr\'echet--Hoeffding bounds; no bibliography entry & proved here: Proposition~\ref{prop:sharp} (upper), Remark~\ref{rem:pairwise} (lower) \\
comonotone sum is maximal in convex order & Meilijson and N\'adas (1979); Dhaene et al.\ (2002) & the convex-order statement is theirs; Lemma~\ref{lem:comon} proves the variance consequence directly for indicators \\
variance bound for a lattice variable & none used & Lemma~\ref{lem:floor}: an integer-valued $X$ with mean $S$ has $\operatorname{Var}(X) \ge f(1-f)$; proved here, with no cited inequality \\
identified set when some areas are unobserved & commonly named the Manski bound; no bibliography entry & proved here: Lemma~\ref{lem:identified}, for the count and for the population share \\
interval for an estimated identified set & Imbens and Manski (2004), a different object & the union-bound form \eqref{eq:coverage} is proved here; it is not their confidence set \\
monotone quantiles, monotone expectation & elementary & proved, Proposition~\ref{prop:monotone} \\
\hline
\end{tabular}
\caption{Classical results used, and what is done with them.  Proposition~
\ref{prop:sharp} and Lemma~\ref{lem:cyclic} assemble the first three rows; we
make no claim that the assembly is itself a new result.}
\label{tab:attribution}
\end{table}
